% !TeX root = ../../Book5.tex
% 纲要标题：### 16.3 根弦与反射
% 纲要位置：工作记录/计划纲要.md，第 4597 行。

\section{根弦与反射}
\label{b5:sec:1603}

相反根向量的括号产生一个 Cartan 方向。把它适当归一化，就会出现已经熟悉的三维 Lie 代数；根弦的长度和反射系数将由它的有限维表示决定。

% 纲要要点已在下文展开。

% 纲要细目（保留原定写作范围）：
%
% * sl_2 子代数
% * 根弦性质
% * 根反射及整数性
% * 用 sl_3 零权升秩解释根空间一维性，把三根根弦写成完整权链，并以矩阵共轭实现根反射
%

\subsection{\texorpdfstring{\(\mathfrak{sl}_2\) 子代数}{sl2 子代数}}
\label{b5:subsec:1603:01}

\begin{theorem}\label{b5:thm:root-sl2}
设 \(\mathfrak g\) 为有限维复半单 Lie 代数，\(\mathfrak h\subseteq\mathfrak g\) 为 Cartan 子代数，\(B\) 为其 Killing 型，\(\alpha\in\mathfrak h^*\) 为根，\(\mathfrak g_{\pm\alpha}\) 为相应根空间。令 \(t_\alpha\in\mathfrak h\) 满足 \(B(t_\alpha,h)=\alpha(h)\) 对所有 \(h\in\mathfrak h\) 成立，并记 \((\alpha,\alpha)=B(t_\alpha,t_\alpha)\)。则 \((\alpha,\alpha)\ne0\)。取
\(e_\alpha\in\mathfrak g_\alpha\)、\(f_\alpha\in\mathfrak g_{-\alpha}\)
满足 \(B(e_\alpha,f_\alpha)=1\)，并令
\[
h_\alpha=\frac{2t_\alpha}{(\alpha,\alpha)},\qquad
E_\alpha=e_\alpha,\qquad
F_\alpha=\frac{2f_\alpha}{(\alpha,\alpha)}.
\]
则 \([h_\alpha,E_\alpha]=2E_\alpha\)、
\([h_\alpha,F_\alpha]=-2F_\alpha\)、
\([E_\alpha,F_\alpha]=h_\alpha\)，三者构成一个
\(\mathfrak{sl}_2(\mathbb C)\) 子代数。每个根空间均为一维，且与 \(\alpha\) 成比例的根只有 \(\alpha,-\alpha\)。
\end{theorem}
\begin{proof}
由前一节的非退化配对可选到 \(e=e_\alpha,f=f_\alpha\)，满足 \([e,f]=t_\alpha\)。若 \(\alpha(t_\alpha)=0\)，则
\(\operatorname{ad}e,\operatorname{ad}f\) 都与半单算子
\(\operatorname{ad}t_\alpha\) 交换。在后者的特征值为 \(c\) 的特征空间上，
\[
[\operatorname{ad}e,\operatorname{ad}f]
=\operatorname{ad}t_\alpha=cI.
\]
取迹得到 \(c\dim V_c=0\)，所以所有特征值为零。
半单性给出 \(\operatorname{ad}t_\alpha=0\)，继而 \(t_\alpha=0\)，与 \(\alpha\ne0\) 矛盾。因此
\(\alpha(t_\alpha)=(\alpha,\alpha)\ne0\)，所列括号关系随即成立。

从原根 \(\alpha\) 出发，若当前根的一半仍是根，就继续减半。由于 \(\Phi\) 有限，这个过程终止于同一直线上的根 \(\gamma\)，满足 \(\gamma/2\notin\Phi\)。前面的计算也适用于 \(\gamma\)，因而可构造相应的三元组 \(E_\gamma,F_\gamma,h_\gamma\)。将这份 \(\mathfrak{sl}_2\) 作用限制到
\[
M_\gamma=\mathfrak h\oplus
\bigoplus_{\beta\in\Phi\cap\mathbb C\gamma}\mathfrak g_\beta.
\]
由\cref{b5:thm:sl2-complete-reducibility}，它是有限维不可约 \(\mathfrak{sl}_2\) 模的直和；其中 \(h_\gamma\) 在 \(\mathfrak g_{c\gamma}\) 上的特征值为 \(2c\)，所以 \(2c\) 为整数。

三维伴随子模
\[
S_\gamma=\mathbb CE_\gamma\oplus\mathbb Ch_\gamma\oplus\mathbb CF_\gamma
\]
同构于不可约模 \(L_2\)。完全可约性允许选取子模补空间 \(W\)，使 \(M_\gamma=S_\gamma\oplus W\)，再把 \(W\) 分解为不可约模。在\cref{b5:prop:sl2-model}中令标准生成元 \(e,f,h\) 分别对应 \(E_\gamma,F_\gamma,h_\gamma\)，逐个代入模型，就可以确定每种直和项对权 \(0\) 到权 \(2\) 的升算子贡献了多少秩。

平凡模 \(L_0\) 只有权 \(0\)，升算子为零，贡献的秩为 \(0\)。若最高权为奇数 \(m=2r+1\)，模型中的 \(v_r\) 有权 \(m-2r=1\)。在 \(M_\gamma\) 中，权 \(1\) 只能来自 \(\mathfrak g_{\gamma/2}\)，因为 \(h_\gamma\) 在 \(\mathfrak g_{c\gamma}\) 上的权为 \(2c\)；这与选择 \(\gamma\) 时已经保证 \(\gamma/2\notin\Phi\) 矛盾，所以奇最高权直和项不能出现。若最高权为偶数 \(m=2r\geq2\)，则 \(v_r\) 的权为 \(0\)，\(v_{r-1}\) 的权为 \(2\)，且
\[
E_\gamma v_r=r(r+1)v_{r-1}\ne0.
\]
两权空间都一维，所以这一项贡献的秩恰为 \(1\)。

另一方面，整个 \(M_\gamma\) 的零权空间是 \(\mathfrak h\)，权 \(2\) 空间是 \(\mathfrak g_\gamma\)，并且升算子在零权空间上为
\[
\mathfrak h\longrightarrow\mathfrak g_\gamma,\qquad
h\longmapsto[E_\gamma,h]=-\gamma(h)E_\gamma.
\]
它的像恰为 \(\mathbb CE_\gamma\)：\(\gamma\ne0\)，且在 \(h_\gamma\) 上的值为 \(-2E_\gamma\)，故总秩为 \(1\)。直和分解中的升算子逐项作用，像也为直和，所以总秩是各项贡献的秩之和。已知的 \(S_\gamma\cong L_2\) 已经贡献了秩 \(1\)，于是 \(W\) 不能再含任何非平凡偶最高权直和项，只能由 \(L_0\) 组成。因此 \(M_\gamma\) 的非零权只有 \(\pm2\)，两权空间都一维，即
\[
\Phi\cap\mathbb C\gamma=\{\gamma,-\gamma\},\qquad
\dim\mathfrak g_\gamma=\dim\mathfrak g_{-\gamma}=1.
\]
原根 \(\alpha\) 与 \(\gamma\) 共线，故 \(\alpha=\pm\gamma\)，从而所证结论也适用于原根 \(\alpha\)。
\end{proof}

两套根向量的归一化各有用途。Killing 配对为 \(1\) 的 \(e_\alpha,f_\alpha\) 适合 Casimir 计算；满足标准 \(\mathfrak{sl}_2\) 关系的则是 \(E_\alpha,F_\alpha,h_\alpha\)。除非恰有 \((\alpha,\alpha)=2\)，不得把两套记号不加说明地混用。

\begin{example}\label{b5:exm:sl3-zero-weight-rank}
在 \(\mathfrak g=\mathfrak{sl}_3(\mathbb C)\) 中，取迹零对角 Cartan 子代数 \(\mathfrak h\)，令
\(\varepsilon_i(\operatorname{diag}(t_1,t_2,t_3))=t_i\)、
\(\alpha=\varepsilon_1-\varepsilon_2\)。相应的标准三元组为
\[
E_\alpha=E_{12},\qquad F_\alpha=E_{21},\qquad
H=h_\alpha=E_{11}-E_{22}.
\]
令 \(K=\operatorname{diag}(1,1,-2)\)，则
\(\mathfrak h=\mathbb CH\oplus\mathbb CK\)，并且
\[
M_\alpha=\mathbb CE_{12}\oplus\mathfrak h\oplus\mathbb CE_{21}.
\]
在有序基 \((E_{12},H,E_{21},K)\) 上，伴随作用的矩阵为
\[
\begin{aligned}
[\operatorname{ad}H]&=\operatorname{diag}(2,0,-2,0),\\
[\operatorname{ad}E_{12}]&=
\begin{pmatrix}0&-2&0&0\\0&0&1&0\\0&0&0&0\\0&0&0&0\end{pmatrix},\\
[\operatorname{ad}E_{21}]&=
\begin{pmatrix}0&0&0&0\\-1&0&0&0\\0&2&0&0\\0&0&0&0\end{pmatrix}.
\end{aligned}
\]
这里前三个基向量生成三维伴随子模 \(L_2\)，而 \(K\) 与三元组的三个元素都交换，生成平凡子模 \(L_0\)。因而
\(M_\alpha\cong L_2\oplus L_0\)，与\cref{b5:exm:sl3-adjoint-root-strings}的分解一致。

关于 \(H\) 的零权空间为整个 \(\mathfrak h\)，权 \(2\) 空间为 \(\mathbb CE_{12}\)。以上述 \((H,K)\) 为定义域的基、\(E_{12}\) 为值域的基，升算子的矩阵恰为
\[
\left[\operatorname{ad}E_{12}:
(M_\alpha)_0\longrightarrow(M_\alpha)_2\right]
=\begin{pmatrix}-2&0\end{pmatrix}.
\]
所以总秩为一：\(L_2\) 的零权方向 \(H\) 已经贡献这一秩，\(L_0\) 的方向 \(K\) 不作贡献。一般证明中，每增加一份非平凡偶最高权模，零权到权 \(2\) 的映射就会多贡献一个秩；既有的 \(L_2\) 已用尽总秩，因此不能再出现这样的直和项。排除奇最高权项则另由 \(\gamma/2\notin\Phi\) 保证。这两步共同控制根空间的一维性，所比较的不是零权空间本身的维数。
\end{example}

\subsection{根弦性质}
\label{b5:subsec:1603:02}

\begin{theorem}[根弦]\label{b5:thm:root-string}
设 \(\mathfrak g\) 为有限维复半单 Lie 代数，\(\mathfrak h\subseteq\mathfrak g\) 为 Cartan 子代数，\(\Phi\) 为相应根集合，\(\mathfrak g_\gamma\) 为根 \(\gamma\) 的根空间。取 \(\alpha,\beta\in\Phi\) 且 \(\beta\ne\pm\alpha\)，记 \(h_\alpha\in\mathfrak h\) 为根 \(\alpha\) 的余根。则存在非负整数 \(p,q\)，使
\[
\Phi\cap(\beta+\mathbb Z\alpha)
=\{\,\beta-p\alpha,\ldots,\beta+q\alpha\,\},
\qquad
\beta(h_\alpha)=p-q.
\]
这列根对应的根空间之和是根 \(\mathfrak{sl}_2\) 子代数的不可约模；相邻根空间之间的升降括号非零。特别地，若 \(\alpha+\beta\in\Phi\)，则
\([\mathfrak g_\alpha,\mathfrak g_\beta]=\mathfrak g_{\alpha+\beta}\)。
\end{theorem}
\begin{proof}
由前一定理，\(\beta\) 与 \(\alpha\) 不成比例，故所有 \(\beta+k\alpha\) 都非零。令
\[
M=\bigoplus_{\substack{k\in\mathbb Z\\\beta+k\alpha\in\Phi}}
\mathfrak g_{\beta+k\alpha}.
\]
根空间的括号关系使它成为有限维根 \(\mathfrak{sl}_2\) 模。其 \(h_\alpha\) 特征值均为
\(\beta(h_\alpha)+2k\)，每个存在的特征空间都一维。由\cref{b5:thm:sl2-complete-reducibility}，\(M\) 完全可约，全部权为整数。特别地，\(\mathfrak g_\beta\ne0\) 给出权 \(\beta(h_\alpha)\)，所以 \(\beta(h_\alpha)\in\mathbb Z\)。于是所有 \(\beta(h_\alpha)+2k\) 都具有同一奇偶性，其不可约直和项的最高权也具有这一奇偶性。而两个偶最高权模都含权 \(0\)，两个奇最高权模都含权 \(1\)，这都会与相应特征空间至多一维矛盾，故 \(M\) 不可能有两个非零直和项。因此 \(M\) 不可约，其权按步长 \(2\) 连续排列，且关于零对称。首、尾权的和为零，给出
\[
\beta(h_\alpha)-2p+\beta(h_\alpha)+2q=0.
\]
这就是 \(p-q=\beta(h_\alpha)\)。不可约 \(\mathfrak{sl}_2\) 模的升降公式还说明，除两端以外，相邻权空间之间的映射均非零。
\end{proof}

例如在 \(\mathfrak{sl}_3\) 中，取
\(\alpha=\varepsilon_1-\varepsilon_2\)、
\(\beta=\varepsilon_2-\varepsilon_3\)。
根弦为 \(\beta,\beta+\alpha\)，故 \(p=0,q=1\)，
\(\beta(h_\alpha)=-1\)。矩阵等式
\([E_{12},E_{23}]=E_{13}\) 正好实现这一步升根。

\begin{corollary}\label{b5:cor:b2-three-root-string}
设 \(\mathfrak g\) 为有限维复半单 Lie 代数，\(\mathfrak h\) 为其 Cartan 子代数，\(\mathfrak h^*\) 有基 \(\varepsilon_1,\varepsilon_2\)，且根集合为
\[
\Phi=\{\,\pm\varepsilon_1,\ \pm\varepsilon_2,
\ \pm(\varepsilon_1+\varepsilon_2),
\ \pm(\varepsilon_1-\varepsilon_2)\,\}.
\]
取 \(\alpha=\varepsilon_2\)、\(\beta=\varepsilon_1-\varepsilon_2\)，并取根 \(\alpha\) 的标准三元组 \((E_\alpha,F_\alpha,h_\alpha)\)。则经过 \(\beta\) 的 \(\alpha\)-根弦为
\[
\varepsilon_1-\varepsilon_2,\qquad
\varepsilon_1,\qquad
\varepsilon_1+\varepsilon_2.
\]
这三个根空间关于 \(h_\alpha\) 的权依次为 \(-2,0,2\)，其和
\[
M=\mathfrak g_{\varepsilon_1-\varepsilon_2}
\oplus\mathfrak g_{\varepsilon_1}
\oplus\mathfrak g_{\varepsilon_1+\varepsilon_2}
\]
在根 \(\mathfrak{sl}_2\) 的伴随作用下同构于 \(L_2\)，相邻根空间之间的两个升映射和两个降映射均非零。
\end{corollary}
\begin{proof}
在所列根集合中，\(\beta+k\alpha=\varepsilon_1+(k-1)\varepsilon_2\) 为根恰好发生于 \(k=0,1,2\)。因 \(\beta\ne\pm\alpha\)，\cref{b5:thm:root-string}适用，给出 \(p=0,q=2\) 以及 \(\beta(h_\alpha)=-2\)。每增加一个 \(\alpha\)，权增加 \(\alpha(h_\alpha)=2\)，所以三个权为 \(-2,0,2\)。同一定理说明 \(M\) 不可约，而每个根空间一维，故它是三维模 \(L_2\)。

还可把全部升降系数写出。取非零
\(v_0\in\mathfrak g_{\varepsilon_1+\varepsilon_2}\)，令
\(v_1=[F_\alpha,v_0]\)、\(v_2=[F_\alpha,v_1]\)。标准模型的公式给出
\[
\begin{aligned}
[h_\alpha,v_j]&=(2-2j)v_j\quad(0\le j\le2),\\
[F_\alpha,v_0]&=v_1,& [F_\alpha,v_1]&=v_2,&
[F_\alpha,v_2]&=0,\\
[E_\alpha,v_0]&=0,& [E_\alpha,v_1]&=2v_0,&
[E_\alpha,v_2]&=2v_1.
\end{aligned}
\]
其中 \(v_1,v_2\ne0\)，分别生成中间和末端的根空间；因此所述四个相邻映射全都非零。
\end{proof}

所列根集合通常记为 \(B_2\)。中间空间 \(\mathfrak g_{\varepsilon_1}\) 的权 \(0\) 是指单个 \(h_\alpha\) 的特征值为零；它对应的根 \(\varepsilon_1\) 本身非零，因而不属于 Cartan 子代数。这与前例中把 \(\mathfrak h\) 纳入 \(M_\alpha\) 的情形不同。根弦定理要求 \(\beta\ne\pm\alpha\)，正是为了排除沿弦经过零泛函而引入整个 Cartan 子代数的情况。

\ProseParagraphBreak

若交换 \(\alpha,\beta\) 的角色，则 \(\alpha+k\beta=k\varepsilon_1+(1-k)\varepsilon_2\) 为根恰好发生于 \(k=0,1\)。这条弦只有 \(\varepsilon_2,\varepsilon_1\)，所以 \(p=0,q=1\)、\(\alpha(h_\beta)=-1\)，两权为 \(-1,1\)，根空间之和为根 \(\beta\) 的三元组下的 \(L_1\)。两个方向的整数值 \(-2,-1\) 不同，分别记录了三根链与两根链。

\subsection{根反射与整数性}
\label{b5:subsec:1603:03}

\begin{proposition}\label{b5:prop:root-reflection}
设 \(\mathfrak g\) 为有限维复半单 Lie 代数，\(\mathfrak h\) 为其 Cartan 子代数，\(\Phi\subseteq\mathfrak h^*\) 为根集合，\((\ ,\ )\) 为 Killing 型诱导的对偶型。对根 \(\alpha\in\Phi\)，取相应的标准 \(\mathfrak{sl}_2\) 三元组 \((E_\alpha,F_\alpha,h_\alpha)\)，其中 \(h_\alpha\in\mathfrak h\) 为余根。则在 \(\mathfrak h^*\) 上的变换
\[
s_\alpha(\lambda)
=\lambda-\lambda(h_\alpha)\alpha
=\lambda-\frac{2(\lambda,\alpha)}{(\alpha,\alpha)}\alpha
\]
保持根集合，且 \(\beta(h_\alpha)\in\mathbb Z\) 对所有根 \(\beta\) 成立。
内自同构
\[
n_\alpha=\exp(\operatorname{ad}E_\alpha)
\exp(-\operatorname{ad}F_\alpha)
\exp(\operatorname{ad}E_\alpha)
\]
保持 \(\mathfrak h\)，在其上作用为
\(h\mapsto h-\alpha(h)h_\alpha\)，在根上实现 \(s_\alpha\)。
\end{proposition}
\begin{proof}
对非比例根，根弦两端的关系说明
\(\beta-(p-q)\alpha\) 仍在同一根弦中，即 \(s_\alpha\beta\) 仍为根。
比例根只有 \(\pm\alpha\)，其整数值为 \(\pm2\)，反射互换二者。

在三维子代数中代入标准括号，三个有限指数依次给出
\[
\begin{aligned}
\exp(\operatorname{ad}E_\alpha)h_\alpha&=h_\alpha-2E_\alpha,\\
\exp(-\operatorname{ad}F_\alpha)(h_\alpha-2E_\alpha)&=-h_\alpha-2E_\alpha,\\
\exp(\operatorname{ad}E_\alpha)(-h_\alpha-2E_\alpha)&=-h_\alpha.
\end{aligned}
\]
所以 \(n_\alpha(h_\alpha)=-h_\alpha\)。而
\(\ker\alpha\subseteq\mathfrak h\) 与 \(E_\alpha,F_\alpha\) 都交换，所以被固定。
把 \(h\) 写成
\(h-\alpha(h)h_\alpha/2+\alpha(h)h_\alpha/2\)，便得到所列公式。
若 \(x\in\mathfrak g_\beta\)，自同构保持括号，因而
\[
[h,n_\alpha x]
=n_\alpha[n_\alpha^{-1}h,x]
=(s_\alpha\beta)(h)n_\alpha x.
\]
所以它把相应根空间同构地送到反射后的根空间。
\end{proof}

\begin{proposition}\label{b5:prop:sl3-root-reflection-matrix}
在 \(\mathfrak g=\mathfrak{sl}_3(\mathbb C)\) 中，令 \(\mathfrak h\) 为迹零对角 Cartan 子代数，
\(\varepsilon_i(\operatorname{diag}(t_1,t_2,t_3))=t_i\)，并取
\(\alpha=\varepsilon_1-\varepsilon_2\)。对标准三元组
\[
E_\alpha=E_{12},\qquad F_\alpha=E_{21},\qquad
h_\alpha=E_{11}-E_{22},
\]
内自同构
\(n_\alpha=\exp(\operatorname{ad}E_\alpha)
\exp(-\operatorname{ad}F_\alpha)
\exp(\operatorname{ad}E_\alpha)\)
就是矩阵 \(N_\alpha\) 的共轭作用，其中
\[
N_\alpha=(I+E_{12})(I-E_{21})(I+E_{12})
=\begin{pmatrix}0&1&0\\-1&0&0\\0&0&1\end{pmatrix}.
\]
对 \(a+b+c=0\)，有
\[
n_\alpha\bigl(\operatorname{diag}(a,b,c)\bigr)
=\operatorname{diag}(b,a,c),\qquad
n_\alpha(E_{23})=E_{13}.
\]
因此它在 Cartan 子代数上交换前两个对角坐标，在根空间上把
\(\mathfrak g_{\varepsilon_2-\varepsilon_3}\) 送到
\(\mathfrak g_{\varepsilon_1-\varepsilon_3}\)，实现根反射 \(s_\alpha\)。
\end{proposition}
\begin{proof}
若 \(A^2=0\)，则在矩阵空间上 \((\operatorname{ad}A)^3=0\)，并且
\[
\begin{aligned}
\exp(\operatorname{ad}A)(X)
&=X+[A,X]+\frac12[A,[A,X]]\\
&=X+AX-XA-AXA=(I+A)X(I-A).
\end{aligned}
\]
由于 \((I+A)^{-1}=I-A\)，这就是 \(I+A\) 的共轭作用。分别取 \(A=E_{12}\) 和 \(A=-E_{21}\)，可知三个指数的复合为 \(N_\alpha\) 的共轭作用。其左上角的二阶矩阵乘积为
\[
\begin{pmatrix}1&1\\0&1\end{pmatrix}
\begin{pmatrix}1&0\\-1&1\end{pmatrix}
\begin{pmatrix}1&1\\0&1\end{pmatrix}
=\begin{pmatrix}0&1\\-1&0\end{pmatrix};
\]
在第三个坐标上三个因子均为恒等，得到所列 \(N_\alpha\)。直接共轭给出
\[
N_\alpha\operatorname{diag}(a,b,c)N_\alpha^{-1}
=\operatorname{diag}(b,a,c)
=\operatorname{diag}(a,b,c)-(a-b)h_\alpha.
\]
它与\cref{b5:prop:root-reflection}在 \(\mathfrak h\) 上的公式相同。若标准列向量为 \(u_1,u_2,u_3\)，则
\(N_\alpha u_1=-u_2\)、\(N_\alpha u_2=u_1\)、\(N_\alpha u_3=u_3\)。据此，\(N_\alpha E_{23}N_\alpha^{-1}=E_{13}\)，它把根 \(\varepsilon_2-\varepsilon_3\) 的方向实际送到根 \(\varepsilon_1-\varepsilon_3\) 的方向；同一坐标交换还把 \(\alpha\) 送到 \(-\alpha\)。
\end{proof}
